% ---
% id: meat-core/lem:memory-lift
% kind: lem
% title: "Approximating an attack by all-at-once attacks with finite memory"
% summary: >-
%   Every attack on a run, its memory of any dimension and its signals taken in any order, is approximated in trace norm by attacks with finite memory that take every signal at the start. A trace-distance bound that holds for those therefore holds for every attack with the same ε.
% prominence: standard
% status: accepted
% uses: [foundations/lem:register-compression, def:postprocessing-model, reductions/def:protocol]
% mentions: [def:roundmodel, thm:keylength, cor:composed-security, reductions/lem:source-transfer, reductions/rem:why-exact]
% assumes: []
% refs:
%   - "tupkary2026decoy, Thm. A.1, Rem. A.2"
%   - "tupkary2026decoy, Def. 8.2"
%   - "tupkary2026decoy, Sec. 8.2"
%   - "tupkary2026decoy, Lemma 8.1"
% provenance: newly-derived
% lean:
% notes: "Deliberately applied to the physical run, before the reductions, where tupkary2026decoy applies its Thm. A.1 to the analysed protocol after the source map and the squasher, under the finite-dimensional Alice and Bob of its Rem. A.2: compressing only the adversary's registers and the signals in her hands needs no finite dimension of the measured or the transmitted system, and the reductions and the accumulation theorem keep their finite-dimensional registers. The release schedule covers every order in which the source may release its signals, k_j = j being one signal per round and k_1 = n the all-at-once type of reductions/lem:source-transfer. The lemma reads no security claim, so the claims that state the bounds can use it."
% ---

The bound of Theorem~\ref{thm:keylength} is proved for the attacks the
\tref{round model} of Definition~\ref{def:roundmodel} represents: their
registers are finite-dimensional, as every register of that model is, and in
\tref{trusted-source mode} their first channel receives every transmitted
signal, the type the reductions transfer
(\xref{reductions/lem:source-transfer}) and the timing the published analysis
moves to \cite[Lemma~8.1]{tupkary2026decoy}. A physical adversary may keep a
memory of any dimension and receives the signals as the source releases
them; where a signal lives in an infinite-dimensional space, a finite memory
cannot hold it for a later round. The lemma closes that gap by a limit, the
route of \cite[Def.~8.2, Sec.~8.2, Thm.~A.1, Rem.~A.2]{tupkary2026decoy}:
compressing the adversary's registers, and each signal released after the
first round, onto finite-dimensional subspaces
(\xref{foundations/lem:register-compression}) turns every attack into one of
the analysed type whose output state is as close as wanted, and a bound on a
trace distance, which no dimension enters, passes to the limit.
Corollary~\ref{cor:composed-security} applies it. The compression touches
neither the \tref{measured system} nor the \tref{transmitted system} in
transit, so no finite dimension of either is needed; that is why the limit is
taken at the physical run rather than at the analysed protocol. Unlike the
transfers of the reductions module, which hold for every criterion because
they are exact (\xref{reductions/rem:why-exact}), the limit needs a criterion
continuous in the output state, and the trace distances of the security
statement are.

\begin{lemma}\label{lem:memory-lift}
Read a run of the protocol and its attacks as follows.
\begin{enumerate}[label=(\alph*),leftmargin=2.4em]
\item \emph{(the run)} Its quantum phase is an $n$-round protocol of
  \xref{reductions/def:protocol}, whose measured system $\Bsys_j$ of round $j$
  is separable and of any dimension, followed by the
  \tref{fixed-length post-processing model} of
  Definition~\ref{def:postprocessing-model} in the
  \tref{honest-authentication setting}. In trusted-source mode its source
  releases, for $i=1,\dots,n$, a separable transmitted system
  $\Asig_i$ of any dimension, prepared together with Alice's letter register
  $X_i$ in a state $\rho_i$ on $X_i\Asig_i$ that depends on nothing else, and
  the run reads $X_i$ from round $i$ on; in \tref{untrusted-source mode} there
  are no $X_i$ and no $\Asig_i$. Every system of the run other than the
  $\Bsys_j$ and the $\Asig_i$ is finite-dimensional, and its public messages
  are classical registers it keeps.
\item \emph{(attacks)} An \emph{attack} consists of a release schedule
  $0=k_0\le k_1\le\dots\le k_n=n$ with $k_j\ge j$ for every $j$, a state
  $\zeta$ on the adversary's initial register $E_0$, and channels, in the
  sense of \xref{foundations/lem:register-compression},
  \begin{equation}\label{eq:lift-attack}
    \Aev_j\in
    \begin{cases}
      \mathrm{CPTP}\big(\Asig_{k_{j-1}+1}\cdots\Asig_{k_j}E_{j-1},\
        \Bsys_jE'_j\big), & \text{trusted-source mode},\\[0.3ex]
      \mathrm{CPTP}\big(E_{j-1},\ A_j\Bsys_jE'_j\big),
        & \text{untrusted-source mode},
    \end{cases}
  \end{equation}
  for $j=1,\dots,n$, where $E'_j$ is the adversary's memory after round $j$,
  $E_j$ is $E'_j$ together with a copy of the public messages of round $j$,
  and $\Hs{E_0}$ and every $\Hs{E'_j}$ are Hilbert spaces of any dimension.
  Under the attack the run starts from $\zeta$ and its own registers, prepares
  $\rho_i$ in round $j$ for $k_{j-1}<i\le k_j$ before $\Aev_j$ acts, takes
  $\Bsys_j$, and in untrusted-source mode Alice's system $A_j$, from the
  output of $\Aev_j$, and places a copy of the public messages of round $j$ in
  $E_j$. Its \emph{output state} is the state of the run's
  \tref{output registers}, its public messages and $E'_n$. The attack is
  \emph{all-at-once} if $k_1=n$, and has \emph{finite registers} if $E_0$ and
  every $E'_j$ are finite-dimensional.
\end{enumerate}
Then:
\begin{enumerate}[label=(\roman*),leftmargin=2.4em]
\item \emph{(approximation)} for every attack and every $\eta>0$ there is an
  all-at-once attack with finite registers whose memory after round $n$ is a
  subspace of $\Hs{E'_n}$, such that the output states $\rho$ and
  $\rho^{\circ}$ of the two attacks satisfy
  $\|\iota(\rho^{\circ})-\rho\|_1\le\eta$, with $\iota$ the inclusion of that
  subspace into $\Hs{E'_n}$, acting as the identity on every other register;
\item \emph{(trace-distance criteria)} let $L_1$ and $L_2$ be completely
  positive maps, not increasing the trace, from the run's output registers
  and public messages to a finite-dimensional system, put $L\coloneqq
  L_1-L_2$, and for a state $\sigma$ of those registers and the adversary's
  memory after round $n$ put
  $f(\sigma)\coloneqq\frac12\|(L\otimes\idmap)(\sigma)\|_1$; if
  $f(\rho^{\circ})\le\varepsilon$ at the output state $\rho^{\circ}$ of every
  all-at-once attack with finite registers, then $f(\rho)\le\varepsilon$ at
  the output state $\rho$ of every attack;
\item \emph{(processing by the adversary)} let, in (ii), each $L_t$ keep the
  public messages and act on the other registers given their value $m$,
  $L_t(\dyad{m}\otimes x)=\dyad{m}\otimes L_{t,m}(x)$, and let
  $\Phi=\sum_m\dyad{m}\otimes\Phi_m$ act on a state that is classical on the
  public messages by a channel $\Phi_m$ on the adversary's register, given
  their value $m$, keeping them. Then $f((\idmap\otimes\Phi)(\sigma))\le
  f(\sigma)$ for every state $\sigma$ classical on the public messages.
\end{enumerate}
\end{lemma}

\begin{proof}
\emph{Idea.} Written as a sequence of channels, the run is compressed at
finitely many points: the adversary's initial register, her memory after each
of her channels, and every signal released after the first round, as soon as
it is released. The compressions cost at most $\eta$ in trace norm
(\xref{foundations/lem:register-compression}\,(ii), (iii)), and the compressed
run is the run under an all-at-once attack with finite registers: the
adversary takes every signal at the start, compresses those she would have
received later and keeps them, compressed, in her memory until their round.
A trace distance of the output state moves by at most the distance of the
states, which gives (ii); a channel on the adversary's side commutes past the
criterion's maps and does not increase a trace norm, which gives (iii).

\emph{Step 1 (the run as a sequence of channels).} Fix an attack. The run's
own registers are finite-dimensional and start in fixed states, which we take
as written by the channels that first write them. Round $j$ is the composite
of three channels: the preparation $\mathcal{Q}_j$, which appends
$\bigotimes_{k_{j-1}<i\le k_j}\rho_i$ (the identity in untrusted-source mode);
$\Aev_j$, extended by the identity on the run's registers; and the round map
$\Grnd_j$ of \xref{reductions/def:protocol}, extended by the identity on
$E'_j$. After round $n$ the run applies $\Kfin$, which includes the
post-processing, extended by the identity on $E'_n$. These extensions are
channels acting on products as the named map tensored with the identity. A
channel extended by the identity on a finite-dimensional register applies it
to every entry of a block matrix, which is a channel by complete positivity.
$\Grnd_j$ is the \tref{measurement channel} of the measured system followed
by processing of finite-dimensional registers; the measurement channel
$\tau\mapsto\sum_v\tr[M_v\tau]\dyad{v}$ extends to
$\tau\mapsto\sum_v\dyad{v}\otimes\tr_{\Bsys_j}\big[(\sqrt{M_v}\otimes\idon{})
\tau(\sqrt{M_v}\otimes\idon{})\big]$, a sum of conjugations followed by
partial traces and the appending of states, and a completely positive map
$\Psi$ between finite-dimensional systems is a finite sum
$\sum_tK_t(\cdot)K_t^{\dagger}$, which extends to
$\sum_t(K_t\otimes\idon{})(\cdot)(K_t\otimes\idon{})^{\dagger}$. The
output state is the result of these channels applied in order to $\zeta$.

\emph{Step 2 (the compressions).} The compression points are: $E_0$ before
round one; $E'_j$ right after $\Aev_j$, for every $j$; and, for every $i>k_1$,
the transmitted system $\Asig_i$ right after the preparation $\mathcal{Q}_j$
that appends $\rho_i$, before $\Aev_j$ receives it, one signal at a time with
identity channels between. There are at most $2n+1$ of them. At each point,
\xref{foundations/lem:register-compression}\,(ii), applied to the run's state
there under the given attack, gives a non-zero finite-dimensional subspace of
the register at which the bound of that part, $2\sqrt{\delta}+\delta$, is at
most $\eta/(2n+1)$; we write $E_0^{\circ}$, $E_j^{\prime\circ}$ and
${\Asig_i}^{\circ}$ for them. Inserting at each point the compression onto its
subspace, and the identity at every other point, part (iii) of that lemma
gives $\|\rho^{\mathcal{C}}-\rho\|_1\le\eta$ for the output
$\rho^{\mathcal{C}}$ of the run so compressed.

\emph{Step 3 (the compressed run under an attack with finite registers).}
Write each compression as $\iota\circ\mathcal{P}$
(\xref{foundations/lem:register-compression}\,(i)). Each inclusion $\iota$ of
$E'_{j}$ commutes with the maps that follow it up to $\Aev_{j+1}$, the round
map $\Grnd_j$, the preparations of round $j+1$ and the compressions of the
signals they release, all of which act on other registers: for the channels
this is part (iv) of that lemma, and for an appended state $\rho$ it is
$\iota(y)\otimes\rho=(\iota\otimes\idmap)(y\otimes\rho)$, the product form of
part (i). Likewise the inclusion of $E_0$ moves to $\Aev_1$, the inclusion of
each ${\Asig_i}^{\circ}$ to the channel $\Aev_j$ that receives $\Asig_i$, and
the inclusion of $E_n^{\prime\circ}$ past $\Grnd_n$ and $\Kfin$ to the end.
Each projection $\mathcal{P}$ stays where it acts: after $\Aev_j$ it joins
it, and after the preparation of $\rho_i$ it gives
$\rho_i^{\circ}\coloneqq(\idmap_{X_i}\otimes\mathcal{P})(\rho_i)$, by the
product form $\mathcal{P}_{R_1R_2}(x\otimes y)=x\otimes\mathcal{P}_{R_2}(y)$
of part (i) with $R_2=X_i$ and $R_1$ the other systems. So
$\rho^{\mathcal{C}}=\iota(\rho^{\mathrm{fs}})$, where $\rho^{\mathrm{fs}}$ is
the output state of the run that prepares $\rho_i$ for $i\le k_1$ and
$\rho_i^{\circ}$ for $i>k_1$, on $X_i{\Asig_i}^{\circ}$, in the round of the
schedule, against the initial state $\mathcal{P}(\zeta)$ on $E_0^{\circ}$ and
the channels
\begin{equation*}
  \Aev^{\circ}_1\coloneqq\mathcal{P}\circ\Aev_1\circ\iota,
  \qquad
  \Aev^{\circ}_j\coloneqq\mathcal{P}\circ\Aev_j\circ\iota\ \ (j\ge2),
\end{equation*}
the projection being that of $E'_j$, and the inclusion that of $E_0$ for
$j=1$ and of $E'_{j-1}$ and every ${\Asig_i}^{\circ}$ received in round $j$ for
$j\ge2$, each acting as the identity on the other registers. Each
$\Aev^{\circ}_j$ is a channel between finite-dimensional adversary's
registers, with the signals of the schedule and $\Bsys_j$, and $A_j$ in
untrusted-source mode, as before.

\emph{Step 4 (an all-at-once attack with the same output).} Put
$E^{\star}_j\coloneqq E_j^{\prime\circ}{\Asig_{k_j+1}}^{\circ}\cdots
{\Asig_n}^{\circ}$, finite-dimensional, so that $E^{\star}_n=E_n^{\prime\circ}$.
The all-at-once attack has the initial state $\mathcal{P}(\zeta)$, the first
channel
\begin{equation*}
  \mathcal{A}^{\star}_1\coloneqq
  \big(\Aev^{\circ}_1\otimes\idmap_{{\Asig_{k_1+1}}^{\circ}\cdots{\Asig_n}^{\circ}}\big)
  \circ\mathcal{P}_{\Asig_{k_1+1}}\circ\dots\circ\mathcal{P}_{\Asig_n}
  \in\mathrm{CPTP}\big(\Asig_1\cdots\Asig_nE_0^{\circ},\ \Bsys_1E^{\star}_1\big),
\end{equation*}
each $\mathcal{P}_{\Asig_i}$ the projection of $\Asig_i$ onto
${\Asig_i}^{\circ}$ with the identity on the other systems, and, for $j\ge2$,
$\mathcal{A}^{\star}_j\coloneqq\Aev^{\circ}_j\otimes
\idmap_{{\Asig_{k_j+1}}^{\circ}\cdots{\Asig_n}^{\circ}}$, from $E^{\star}_{j-1}$
with the public messages of round $j-1$ to $\Bsys_jE^{\star}_j$; the input
of $\Aev^{\circ}_j$ is a factor of that register. In untrusted-source mode it
is the attack of Step~3 itself. Its run prepares
$\rho_1\otimes\dots\otimes\rho_n$ at the start, and $\mathcal{A}^{\star}_1$
turns each factor $\rho_i$ with $i>k_1$ into $\rho_i^{\circ}$, by the product
form of part (i). From then on no map acts on $X_i{\Asig_i}^{\circ}$ before
the round that reads $X_i$, which is not earlier than the round of the
schedule since $k_i\ge i$, or before the channel that receives
${\Asig_i}^{\circ}$; every map before them acts on other registers, and a map
$\Psi$ extended by the identity on the finite-dimensional
$X_i{\Asig_i}^{\circ}$ takes $y\otimes\rho_i^{\circ}$ to
$\Psi(y)\otimes\rho_i^{\circ}$. Appending $\rho_i^{\circ}$ at the start or in
its round of the schedule therefore leaves every later state the same, and
the output state $\rho^{\circ}$ of the all-at-once attack is
$\rho^{\mathrm{fs}}$. With Steps~2 and~3,
$\|\iota(\rho^{\circ})-\rho\|_1=\|\rho^{\mathcal{C}}-\rho\|_1\le\eta$, which
is (i).

\emph{Step 5 (part (ii)).} The maps $L_1\otimes\idmap$ and $L_2\otimes\idmap$
are completely positive and do not increase the trace, by the Kraus form of
Step~1, so each contracts the trace norm of a Hermitian operator: for
$Z=Z^{+}-Z^{-}$, $\|(L_t\otimes\idmap)(Z)\|_1\le\tr(L_t\otimes\idmap)(Z^{+})
+\tr(L_t\otimes\idmap)(Z^{-})\le\|Z\|_1$. Hence for two states $\sigma$,
$\sigma'$,
\begin{equation*}
  |f(\sigma)-f(\sigma')|\le\tfrac12\big\|(L\otimes\idmap)(\sigma-\sigma')\big\|_1
  \le\tfrac12\sum_{t=1,2}\big\|(L_t\otimes\idmap)(\sigma-\sigma')\big\|_1
  \le\|\sigma-\sigma'\|_1 .
\end{equation*}
Given an attack and $\eta>0$, take the all-at-once attack of (i). By part
(iv) of \xref{foundations/lem:register-compression}, with $G=L$,
$(L\otimes\idmap_{E'_n})\circ\iota=\iota\circ(L\otimes\idmap_{E_n^{\prime\circ}})$,
and $\iota$ preserves the trace norm, so $f(\iota(\rho^{\circ}))=f(\rho^{\circ})
\le\varepsilon$. So $f(\rho)\le\varepsilon+\eta$ for every $\eta>0$.

\emph{Step 6 (part (iii)).} Write $\sigma=\sum_m\dyad{m}\otimes\sigma_m$, with
$\sigma_m$ positive on the other registers and the adversary's. Then
\begin{equation*}
  (L_t\otimes\idmap)\big((\idmap\otimes\Phi)(\sigma)\big)
  =\sum_m\dyad{m}\otimes(L_{t,m}\otimes\Phi_m)(\sigma_m)
  =\Phi'\big((L_t\otimes\idmap)(\sigma)\big),
\end{equation*}
with $\Phi'\coloneqq\sum_m\dyad{m}\otimes\Phi_m$ acting on the output of
$L_t$, which keeps the public messages: both sides apply $L_{t,m}$ to the run's
registers and $\Phi_m$ to the adversary's in each block, which commute by
part (iv) of \xref{foundations/lem:register-compression}. So
$(L\otimes\idmap)(\idmap\otimes\Phi)(\sigma)=\Phi'((L\otimes\idmap)(\sigma))$,
and $\Phi'$ is a positive trace-preserving map on block-diagonal operators,
which contracts the trace norm of the Hermitian $(L\otimes\idmap)(\sigma)$ as
in Step~5. Halving gives the claim.
\end{proof}
